\section{Reacción en cadena en el panorama de la capacidad de cómputo}

\subsection{El reajuste del triángulo de la capacidad de cómputo}

Jiang Daxin (蒋大昕) analizó el efecto en cadena que el paradigma de RL tiene sobre el triángulo «algoritmo-capacidad de cómputo-datos»:

\begin{importantbox}{Tres juicios de certeza decreciente}
\begin{enumerate}
\item \textbf{Cierto}: la demanda de capacidad de cómputo del lado de la inferencia crece de manera multiplicativa (test-time scaling)
\item \textbf{Muy probablemente cierto}: la capacidad de cómputo de la etapa de entrenamiento con RL no es menor que la del preentrenamiento; los datos de Self-play no tienen, en teoría, un límite superior (por ejemplo, para entrenar Strawberry se usaron decenas de miles de tarjetas H100 durante varios meses)
\item \textbf{Incierto}: si la cantidad de parámetros del modelo principal necesita seguir haciendo scale; si el RL produce un efecto de «amplificador», el ROI del scaling de parámetros podría volver a ser positivo
\end{enumerate}
\end{importantbox}

\begin{warningbox}{Si se cumple el tercer punto}
Si el efecto de amplificación del RL vuelve a hacer eficaz el scaling de parámetros, entonces el crecimiento de la capacidad de cómputo regresaría a la dimensión cuadrática de $\text{cantidad de parámetros} \times \text{cantidad de datos}$: la demanda sobre toda la infraestructura de cómputo sería explosiva.
\end{warningbox}

\subsection{¿Quién sentirá primero la reevaluación de la capacidad de cómputo?}
Aunque los tres panelistas parten de ángulos distintos, sus juicios llegan a conclusiones cercanas: la presión sobre la capacidad de cómputo en el futuro no se manifestará solo en los clústeres de entrenamiento, sino que se transmitirá cuesta arriba a través de los chips, las plataformas en la nube, los servicios de inferencia y las empresas de aplicaciones. Es decir, lo que trae el RL Scaling no es un aumento de costo puntual, sino que toda la cadena de la industria tiene que recalcular el ROI.

\begin{table}[H]
\centering
\begin{tabular}{p{3.2cm}p{5.2cm}p{5cm}}
\toprule
Rol & Presión directa & Posible respuesta \\
\midrule
Empresas de modelos base & El entrenamiento y el test-time scaling elevan a la vez el presupuesto de GPU & Asignar con más precisión el presupuesto de preentrenamiento, posentrenamiento e inferencia \\
Proveedores de infraestructura en la nube & La demanda pico aumenta y la programación de los clústeres se vuelve más compleja & Ampliar el suministro de tarjetas de gama alta y mejorar el aprovechamiento multiinquilino \\
Empresas emergentes de aplicaciones & El costo de inferencia sube al llamar a modelos más potentes & Reducir el costo por solicitud mediante destilación, caché y flujos de trabajo especializados \\
\bottomrule
\end{tabular}
\caption{El RL Scaling afecta de manera distinta a los diferentes roles de la cadena de la industria}
\end{table}

\begin{knowledgebox}{Un juicio comercial implícito en la mesa redonda}
Si la demanda de capacidad de cómputo del lado de la inferencia estalla primero, quienes se beneficien primero no serán necesariamente los equipos con el modelo de más parámetros, sino posiblemente las empresas \textbf{más hábiles para «empaquetar la inferencia costosa como un servicio de alto valor»}. Por ello, el diseño de producto, la estrategia de caché y la división de flujos de trabajo serán tan importantes como la capacidad del modelo.
\end{knowledgebox}

\subsection{Resumen de la sección}
El cambio en el panorama de la capacidad de cómputo no es simplemente que «las tarjetas se encarecieron», sino que el RL Scaling vuelve a elevar tanto el entrenamiento como la inferencia, y transmite la presión hacia la nube, las aplicaciones y los modelos de negocio. Quien logre convertir una alta capacidad de cómputo en una entrega de alto valor tendrá más probabilidades de tomar la delantera en la siguiente ronda de competencia.

